% RASCUNHO v1 (28/09/2026) -- v0 enviado ao autor (commit 7793f74); v1 acrescenta a secao heuristica HKO. Compilar com ~/.local/bin/tectonic -X compile nota.tex.
% Todos os numeros vem de results/momentos_negativos_20260927/ (analise.json, analise_c8.json, analise_mn3.json).
\documentclass[11pt]{amsart}
\usepackage[margin=1in]{geometry}
\usepackage{amsmath,amssymb,booktabs}
\usepackage[hidelinks]{hyperref}

\newcommand{\Lp}{L'}
\DeclareMathOperator{\med}{mean}

\title[Numerical evidence for negative discrete moments of Dirichlet $L$-functions]{Numerical evidence for conjectures on negative discrete second moments of Dirichlet $L$-functions}
\author{Dionatan Resende}
\address{Independent researcher}
\email{dionatandieresende@gmail.com}
\date{Draft v1, 28 September 2026}

\begin{document}

\begin{abstract}
Pearce-Crump recently conjectured that, for a primitive Dirichlet character $\chi$ modulo $q$,
$\sum_{0<\gamma\le T}|\Lp(\rho,\chi)|^{-2}\sim \frac{A_q}{2\pi}T$ with
$A_q=\frac{6}{\pi^2}\prod_{p\mid q}\frac{p}{p+1}$, and
$\sum_{0<\gamma\le T}|L(2\rho,\chi^2)/\Lp(\rho,\chi)|^{2}\sim\frac{\varphi(q)}{q}\frac{T}{2\pi}$.
These conjectures come from a long-mollifier heuristic and generalise conjectures of Gonek and Ng for the Riemann
zeta-function. We test them numerically for the real primitive characters of conductor $4$, $5$, $8$ and $12$, using
between $2000$ and $3000$ zeros of each $L$-function, computed and verified for completeness, together with the first
$27000$ zeros of $\zeta(s)$. To remove the common finite-height bias, we compare each $L$-function with $\zeta(s)$ at
the same local density of zeros. In all eight tests the predicted arithmetic factor lies inside the $95\%$ bootstrap
interval and the value without arithmetic factor is excluded. For conductor $12$, where the two conjectures predict
$1/2$ and $1/3$, we measure $0.501\,[0.445,0.555]$ and $0.340\,[0.301,0.381]$. All predictions and analysis choices
were fixed before the corresponding computations.
\end{abstract}

\maketitle

\section{Introduction}

Let $\rho=\tfrac12+i\gamma$ denote the non-trivial zeros of $\zeta(s)$. Assuming the Riemann Hypothesis and that all
zeros are simple, Gonek~\cite{Gonek} conjectured
\begin{equation}\label{eq:gonek}
\sum_{0<\gamma\le T}\frac{1}{|\zeta'(\rho)|^{2}}\sim\frac{3}{\pi^{3}}\,T,
\end{equation}
and Milinovich and Ng~\cite{MN} proved a lower bound of half this size. Ng~\cite{Ng} conjectured the companion asymptotic
$\sum_{0<\gamma\le T}|\zeta(2\rho)/\zeta'(\rho)|^{2}\sim T/2\pi$, with a matching lower bound of half the size due to
Sinha~\cite{Sinha}. Negative discrete moments of this kind enter explicit formulae for the summatory functions of the
M\"obius and Liouville functions, and their sizes were modelled with random matrix theory by Hughes, Keating and
O'Connell~\cite{HKO}. Numerically, Hiary and Odlyzko~\cite{HO} found agreement with \eqref{eq:gonek} using large sets
of zeros near the $10^{16}$th, $10^{20}$th and $10^{23}$rd zeros.

For a primitive Dirichlet character $\chi$ modulo $q>1$, Pearce-Crump~\cite{PC} proved lower bounds for
\[
J_\chi(T)=\sum_{0<\gamma\le T}\frac{1}{|\Lp(\rho,\chi)|^{2}},\qquad
K_\chi(T)=\sum_{0<\gamma\le T}\left|\frac{L(2\rho,\chi^{2})}{\Lp(\rho,\chi)}\right|^{2},
\]
uniformly in the conductor, and formulated the following conjectures~\cite[Conjectures~7 and~8]{PC}:
\begin{align}
J_\chi(T)&\sim\frac{A_q}{2\pi}\,T,\qquad A_q=\frac{6}{\pi^{2}}\prod_{p\mid q}\frac{p}{p+1},\label{eq:c7}\\
K_\chi(T)&\sim\frac{\varphi(q)}{q}\,\frac{T}{2\pi}.\label{eq:c8}
\end{align}
Both reduce to the conjectures of Gonek and Ng when $q=1$. They were obtained from Farmer's long-mollifier
heuristic~\cite{Farmer}, which, as Pearce-Crump points out, is not a proof and is applied there to discrete sums
rather than to the continuous mean values for which it was designed. Their derivation from the ratios
conjecture~\cite{CFZ,CS} is left open in~\cite{PC}.

The content of \eqref{eq:c7} and \eqref{eq:c8} beyond the case $q=1$ lies in the arithmetic factors
$\prod_{p\mid q}p/(p+1)$ and $\varphi(q)/q$. The aim of this note is to test these factors numerically. We are not aware
of previous numerical tests of \eqref{eq:c7} or \eqref{eq:c8}; our literature search is described in
Section~\ref{sec:search}.

\section{Method}

\subsection{Normalisation}
Since the zeros of $L(s,\chi)$ at height $t$ have density $\frac{1}{2\pi}\log\frac{qt}{2\pi}$, \eqref{eq:c7} and
\eqref{eq:c8} say, to first order, that the zero-wise averages of
\[
v_n=\frac{\log(q\gamma_n/2\pi)}{|\Lp(\rho_n,\chi)|^{2}},\qquad
w_n=\frac{\log(q\gamma_n/2\pi)\,|L(2\rho_n,\chi^{2})|^{2}}{|\Lp(\rho_n,\chi)|^{2}}
\]
tend to $A_q$ and $\varphi(q)/q$ respectively (to $6/\pi^2$ and $1$ for $\zeta$).

\subsection{A differential statistic}
At the heights accessible to us the convergence is slow: for $\zeta(s)$ the average of $v_n$ over the zeros
$1$--$3000$, $3001$--$10000$ and $10001$--$20000$ equals $0.912$, $0.963$ and $0.992$ times $6/\pi^2$. We therefore
do not compare absolute averages with the conjectured constants. Instead, we compare each $L$-function with
$\zeta(s)$ at the same local density $D=\log(q\gamma/2\pi)$. We split the zeros of $L(s,\chi)$ into ten bins at the
deciles of $D$ and define
\[
R=\frac{\med_{\chi}(v)}{\sum_{b}\frac{N_b}{N}\,\med_{\zeta,\,b}(v)},
\]
where $N_b$ is the number of zeros of $L(s,\chi)$ in bin $b$ and $\med_{\zeta,b}(v)$ is the average of $v$ over the
zeros of $\zeta$ whose density $\log(\gamma/2\pi)$ lies in the same bin; $R_8$ is defined in the same way with $w$. If
the finite-height corrections depend on the height mainly through the density, they cancel in $R$ to first order, and
\eqref{eq:c7}, \eqref{eq:c8} predict
\[
R=\prod_{p\mid q}\frac{p}{p+1},\qquad R_8=\frac{\varphi(q)}{q}.
\]
This cancellation is an assumption, not a theorem; it is supported by the fact that the absolute ratios
$\med(v)/A_q$ for $\chi$ ($0.932$ for $q=4$, $0.937$ for $q=5$) lie in the same range as for $\zeta$ at comparable
heights.

\subsection{Alternatives}
For each test we also record two alternatives: $H_0$, no arithmetic factor ($R=1$); and, for \eqref{eq:c7}, the
factor $\prod_{p\mid q}p^2/(p^2-1)$ obtained by simply removing the Euler factors at $p\mid q$ from $1/\zeta(2)$; for
\eqref{eq:c8}, the factor of \eqref{eq:c7}, which tests whether the two moments require different factors.

\subsection{Uncertainty}
The terms $1/|\Lp(\rho,\chi)|^{2}$ have a heavy tail produced by close pairs of zeros: the largest $1\%$ of the terms
carry $12$--$23\%$ of the sums. We use $95\%$ percentile intervals from a block bootstrap (blocks of $100$ consecutive
zeros, resampled independently for $L(s,\chi)$ and $\zeta(s)$, $2000$ replicates). Because the bootstrap may
under-cover with heavy tails, we also recompute $R$ after removing the five largest terms of $L(s,\chi)$.

\subsection{Pre-registration}
For each group of tests, the predictions, alternatives, statistics, bootstrap settings and decision rules were written
down and fixed (with SHA-256 hashes) before the derivatives were computed. A test is counted as favourable if the
predicted value lies in the interval and $H_0$ does not.

\section{Data}

\subsection{Zeros of $\zeta(s)$}
We use the first $27000$ zeros ($\gamma\le 23476.2$) from Odlyzko's table \texttt{zeros1} (declared accuracy
$3\cdot10^{-9}$). The tests for $q=4,5$ use the first $20000$ zeros, and those for $q=8,12$ use $27000$, so that the
densities of the reference cover those of the $L$-functions.

\subsection{Zeros of $L(s,\chi)$}
We use the real primitive characters $\chi_{-4}=(\frac{-4}{\cdot})$ (odd), $\chi_5=(\frac{\cdot}{5})$,
$\chi_8=(\frac{8}{\cdot})$ and $\chi_{12}=(\frac{12}{\cdot})$ (even). Zeros were located by sign changes of the real
function $Z_\chi(t)=e^{i\theta_\chi(t)}L(\tfrac12+it,\chi)$, with $L$ evaluated by Hurwitz zeta sums in
\texttt{mpmath} at $30$ digits, a scan step of at most $1/8$ of the mean spacing, and refinement to $10^{-10}$. Each
list was checked for (V1) agreement of the count with $\theta_\chi(T)/\pi$ at the end points, (V2) agreement with an
independent rescan at half the step, (V3) a second evaluation at random zeros, and (V4) reality of $Z_\chi$.

Sign-change scans lose pairs of very close zeros. This happened: the original lists for $\chi_{-4}$ and $\chi_5$
failed V1/V2 (deficits of $14$ and $6$ zeros in $12000$), and the list for $\chi_8$ failed V2 ($2000$ found against $2002$
in the rescan). In every case the missing zeros were recovered by rescanning all gaps larger than $1.6$ mean spacings
with a step of $1/32$ of the mean spacing; they were always close pairs (for $\chi_8$, the pair $1470.2269$,
$1470.2791$). The corrected lists pass all checks, and the corrected $\chi_8$ list has exactly the $2002$ zeros found by
the rescan. We use the zeros $1$--$3000$ for $q=4,5$ ($\gamma\le 2891.3$ and $2808.1$) and $1$--$2000$ for $q=8,12$
($\gamma\le 1856.6$ and $1764.1$).

\subsection{Derivatives and $L(2\rho,\chi^{2})$}
Since $|Z_\chi'(\gamma)|=|\Lp(\rho,\chi)|$ at a zero, we computed $|\zeta'(\rho)|$ with \texttt{mpmath.siegelz} and
$|\Lp(\rho,\chi)|$ by a central difference of $Z_\chi$ with step $10^{-6}$ at $20$ digits (relative error below
$4\cdot10^{-10}$ against step $10^{-5}$ at $30$ digits). For real $\chi$, $\chi^{2}$ is the principal character, so
$L(2\rho,\chi^{2})=\zeta(1+2i\gamma)\prod_{p\mid q}(1-p^{-1-2i\gamma})$.

\section{Results}

\begin{table}[h]
\centering
\begin{tabular}{lccccc}
\toprule
$q$ & moment & $R$ measured & $95\%$ interval & predicted & alternatives excluded\\
\midrule
$4$  & $J$ & $0.655$ & $[0.592,0.715]$ & $2/3$ & $1$, $4/3$\\
$5$  & $J$ & $0.827$ & $[0.753,0.901]$ & $5/6$ & $1$, $25/24$\\
$8$  & $J$ & $0.687$ & $[0.605,0.784]$ & $2/3$ & $1$, $4/3$\\
$12$ & $J$ & $0.501$ & $[0.445,0.555]$ & $1/2$ & $1$, $3/2$\\
\midrule
$4$  & $K$ & $0.490$ & $[0.430,0.549]$ & $1/2$ & $1$, $2/3$\\
$5$  & $K$ & $0.753$ & $[0.673,0.835]$ & $4/5$ & $1$ (not $5/6$)\\
$8$  & $K$ & $0.558$ & $[0.450,0.684]$ & $1/2$ & $1$ (not $2/3$)\\
$12$ & $K$ & $0.340$ & $[0.301,0.381]$ & $1/3$ & $1$, $1/2$\\
\bottomrule
\end{tabular}
\caption{Differential statistics $R$ (for $J_\chi$) and $R_8$ (for $K_\chi$). Removing the five largest terms gives
$0.633$ ($q=4$) and $0.794$ ($q=5$) for $J$, without changing any conclusion.}
\label{tab:results}
\end{table}

In all eight tests the predicted factor lies inside the interval and $H_0$ is excluded (Table~\ref{tab:results}). In
six of them the other alternative is excluded as well. In particular, for $q=4$ and $q=12$ the two moments require
their own factors: exchanging the factors of \eqref{eq:c7} and \eqref{eq:c8} is excluded. The characters $\chi_8$ and
$\chi_{-4}$ have the same predicted factors but opposite parity, and their measured values are compatible.

For $\zeta(s)$, the average of $w_n$ (Ng's moment, conjectured limit $1$) equals $0.927$, $0.948$ and $0.934$ on the
zeros $1$--$3000$, $3001$--$10000$ and $10001$--$20000$, again showing finite-height deviations of a few per cent
that the differential statistic is designed to cancel.

\section{A heuristic check of the arithmetic factors}\label{sec:hko}

The constants in \eqref{eq:c7} and \eqref{eq:c8} can also be recovered from the random matrix model of Hughes, Keating
and O'Connell~\cite{HKO}, independently of the mollifier heuristic. For $\zeta(s)$ that model predicts
$\sum_{0<\gamma\le T}|\zeta'(\rho)|^{2k}\sim\frac{T}{2\pi}\frac{G^{2}(k+2)}{G(2k+3)}a(k)(\log T)^{k(k+2)+1}$, where $G$ is
Barnes' function and $a(k)=\prod_p(1-\frac1p)^{k^{2}}\,\mathbb{E}_\theta|1-p^{-1/2}e^{i\theta}|^{-2k}$; for $k=-1$ this
gives \eqref{eq:gonek}. For a primitive character $\chi$ modulo $q$ we keep the same random matrix factor and replace
the local average by $\mathbb{E}_\theta|1-\chi(p)p^{-1/2}e^{i\theta}|^{-2k}$, which equals $1$ when $p\mid q$, while the
factor $(1-\frac1p)^{k^{2}}$ is kept for every prime. For $k=1$ this reproduces the factor $\varphi(q)/q$ in the mean
square $\int_0^T|L(\tfrac12+it,\chi)|^{2}\,dt$, as follows from the approximate functional equation and the mean value
theorem. For $k=-1$ it gives
\[
\prod_{p\nmid q}\Bigl(1-\frac1{p^{2}}\Bigr)\prod_{p\mid q}\Bigl(1-\frac1p\Bigr)=\frac{6}{\pi^{2}}\prod_{p\mid q}\frac{p}{p+1}=A_q .
\]
In the same way, the local average of $|L_p(2\rho,\chi^{2})|^{2}|L_p(\rho,\chi)|^{-2}$ equals $1/(1-\frac1p)$ for
$p\nmid q$ and $1$ for $p\mid q$, which gives the factor $\varphi(q)/q$ of \eqref{eq:c8}. Thus two independent heuristics
lead to the same constants. The alternatives excluded by the data also have a simple meaning in this picture: the
factor $\prod_{p\mid q}p^{2}/(p^{2}-1)$ corresponds to dropping $(1-\frac1p)$ at $p\mid q$, and exchanging the factors
of \eqref{eq:c7} and \eqref{eq:c8} corresponds to using the local average of one moment for the other. This is a
heuristic consistency check and says nothing about lower-order terms; we do not claim it is new.

\section{Discussion and limitations}

The data support the arithmetic factors in \eqref{eq:c7} and \eqref{eq:c8} for four real characters, including a
composite conductor. The evidence is limited in several ways: the heights are low ($\gamma<3000$); only real characters
of small conductor were used; the reference set of $\zeta$-zeros is shared between tests; the bootstrap intervals may
be optimistic because of heavy tails; and the differential statistic relies on the assumption that finite-height
corrections depend on the height mainly through the density of zeros.

A derivation of \eqref{eq:c7} and \eqref{eq:c8} from the ratios conjecture, as suggested in~\cite{PC}, would also
predict the lower-order terms. At the low heights used here those terms are of a few per cent and measurable, so the
present data could test such a prediction in absolute rather than differential form. Complex characters and larger
heights are natural further tests.

\section{Literature search}\label{sec:search}
We searched zbMATH Open, the arXiv, and the works citing Gonek's paper on negative moments (via OpenAlex) for numerical
computations of $\Lp(\rho,\chi)$ at zeros of Dirichlet $L$-functions or of the sums $J_\chi$, $K_\chi$, and found none;
the only numerical study of the corresponding moments of $\zeta$ that we found is~\cite{HO}. MathSciNet was not
available to us. % TODO: completar se houver resposta do autor ou nova busca.

\section*{Data and code}
The zeros, derivatives, pre-registration documents and analysis code are available from the author on request.
% TODO: substituir pelo DOI do Zenodo quando houver deposito.
All computations were run on a single laptop with two cores; the computation and verification of the zeros of the
$L$-functions dominated the cost.

\section*{Acknowledgements}
The computations, analysis and drafting were carried out with substantial assistance from an AI system (Claude,
Anthropic); the author checked the results and takes full responsibility for the content.
% TODO(autor): revisar esta declaracao conforme a politica do arXiv e da revista escolhida.

\begin{thebibliography}{99}
\bibitem{CFZ} J. B. Conrey, D. W. Farmer, M. R. Zirnbauer, Autocorrelation of ratios of $L$-functions, Comm. Number
Theory Phys. 2 (2008), no. 3, 593--636.
\bibitem{CS} J. B. Conrey, N. C. Snaith, Applications of the $L$-functions ratios conjectures, Proc. London Math.
Soc. (3) 94 (2007), no. 3, 594--646.
\bibitem{Farmer} D. W. Farmer, Long mollifiers of the Riemann zeta-function, Mathematika 40 (1993), no. 1, 71--87.
\bibitem{Gonek} S. M. Gonek, The second moment of the reciprocal of the Riemann zeta-function and its derivative, talk
at the Mathematical Sciences Research Institute, Berkeley, June 1999.
\bibitem{HO} G. A. Hiary, A. M. Odlyzko, Numerical study of the derivative of the Riemann zeta function at zeros,
Comment. Math. Univ. St. Pauli 60 (2011), 47--60.
\bibitem{HKO} C. P. Hughes, J. P. Keating, N. O'Connell, Random matrix theory and the derivative of the Riemann zeta
function, Proc. R. Soc. Lond. A 456 (2000), 2611--2627.
\bibitem{MN} M. B. Milinovich, N. Ng, A note on a conjecture of Gonek, Funct. Approx. Comment. Math. 46.2 (2012),
177--187.
\bibitem{Ng} N. Ng, Limiting distributions and zeros of Artin $L$-functions, PhD thesis, University of British Columbia,
2000.
\bibitem{PC} A. Pearce-Crump, Negative discrete second moments of Dirichlet $L$-functions, arXiv:2606.25094 (2026).
\bibitem{Sinha} S. Sinha, Some Conditional Results involving Arithmetic Functions, PhD thesis, University of Missouri,
2025.
\end{thebibliography}

\end{document}
